\documentclass[../../main.tex]{subfiles}% 注意这里的文件路径不能用 ./main.tex ,否则用latexmk编译子文件会报错
\graphicspath{{\subfix{./image/}}{\subfix{../../image/}}} % 指定图片引用路径,后续可以直接使用图片文件名
% 如果图片存放在上级目录,则使用 ../../image/ 作为备用路径

% 例如:
% \begin{figure}[H]
% \centering
% \includegraphics[width=0.3\textwidth]{图.png}
% \caption{}
% \label{figure:图}
% \end{figure}
% 注意:上述\label{}一定要放在\caption{}之后,否则引用图片序号会只会显示??.

\begin{document}

\section{一般测度空间中的一些常用结论}

\begin{theorem}[简单函数逼近定理]\label{theorem:实分析-简单函数逼近定理}
设 $(X,\mathscr{M},\mu)$ 为测度空间，$f:X\to[0,\infty]$ 为非负可测函数。则存在一列非负简单函数 $\{\varphi_n\}_{n=1}^{\infty}$，使得对每个 $x\in X$，有$\varphi_n(x)\leqslant \varphi_{n+1}(x)$ 对所有$n$成立，且 $\lim\limits_{n\to\infty}\varphi_n(x)=f(x)$.
\end{theorem}

\begin{theorem}[Beppo Levi单调收敛定理]\label{theorem:实分析-Beppo Levi单调收敛定理}
设 $(X,\mathscr{M},\mu)$ 为测度空间，$\{f_n\}_{n=1}^{\infty}$ 为一列非负可测函数 $f_n:X\to[0,\infty]$，满足对每个 $x\in X$，有
\begin{align*}
f_n\left( x \right) \leqslant f_{n+1}\left( x \right) \text{对所有}n\text{成立},\text{且}\lim_{n\rightarrow \infty} f_n\left( x \right) =f\left( x \right) ,
\end{align*}
则$f:X\to[0,\infty]$是可测函数且
\[
\lim_{n\to\infty}\int_X f_n\,\mathrm{d}\mu = \int_X \lim_{n\to\infty} f_n\,\mathrm{d}\mu = \int_X f\,\mathrm{d}\mu.
\]
\end{theorem}

\begin{theorem}[Chebyshev不等式]\label{theorem:实分析-Chebyshev不等式}
设 $(X,\mathscr{M},\mu)$ 为测度空间，$0<p<\infty,f:X\to[0,\infty]$ 为非负可测函数。则对$\forall t>0$，有
\begin{align*}
\mu(\{x: |f(x)|\geqslant t\})\leqslant \frac{1}{t^p}\int_X |f(x)|^p\,\mathrm{d}\mu(x).
\end{align*}
\end{theorem}
\begin{proof}
若 $\int_X |f|^p\,\mathrm{d}\mu=+\infty$，则不等式平凡,下设$\int_X |f|^p\,\mathrm{d}\mu<\infty$.由积分的单调性可得
\begin{align*}
&t^p\mu \left( \left\{ x:f(x)\geqslant t \right\} \right) =\int_{\left\{ x:f(x)\geqslant t \right\}}{t^p\,\mathrm{d}\mu (x)}\leqslant \int_{\left\{ x:f(x)\geqslant t \right\}}{\left| f(x) \right|^p\,\mathrm{d}\mu (x)}\leqslant \int_X{\left| f(x) \right|^p\,\mathrm{d}\mu (x)}\\
&\Longrightarrow \mu \left( \left\{ x:f(x)\geqslant t \right\} \right) \leqslant \frac{1}{t^p}\int_X{\left| f(x) \right|^p\,\mathrm{d}\mu (x)}.
\end{align*}
\end{proof}

\begin{theorem}[Tonelli定理]
\label{theorem:Tonelli定理}
\begin{enumerate}[(1)]
\item 设 $(X,\mathscr{M},\mu)$ 和 $(Y,\mathscr{N},\nu)$ 为 $\sigma$--有限的测度空间，$f: X \times Y \to [0,\infty]$ 为乘积空间上的非负可测函数。则
\[
\int_X \left( \int_Y f(x,y)\,\mathrm{d}\nu(y) \right) \mathrm{d}\mu(x)
=
\int_Y \left( \int_X f(x,y)\,\mathrm{d}\mu(x) \right) \mathrm{d}\nu(y)
=
\int_{X\times Y} f\,\mathrm{d}(\mu\times\nu).
\]

\item 设 $(X,\mathscr{M},\mu)$ 和 $(Y,\mathscr{N},\nu)$  为测度空间，
$f:X\times Y\to[0,\infty]$ 为非负可测函数。
若存在可数个
$$
\{A_n\}_{n=1}^{\infty}\subset \mathscr{M},\quad \{B_n\}_{n=1}^{\infty}\subset \mathscr{N},
$$
且满足
\[
\operatorname{supp} f \subset \bigcup_{n=1}^{\infty} (A_n\times B_n),\quad \mu(A_n)<\infty,\nu(B_n)<\infty,\,\, \forall n\in \mathbb{N}.
\]
则
\[
\int_X \left(\int_Y f(x,y)\,\mathrm{d}\nu(y)\right)\mathrm{d}\mu(x)
=
\int_Y \left(\int_X f(x,y)\,\mathrm{d}\mu(x)\right)\mathrm{d}\nu(y)
=
\int_{X\times Y} f\,\mathrm{d}(\mu\times\nu).
\]
特别地,若存在$\sigma$--有限集合$A\subset X,B\subset Y$,满足
\begin{align*}
\operatorname{supp} f \subset A\times B,\quad \mu(A)<\infty ,\nu(B)<\infty.
\end{align*}
则也有
\[
\int_X \left(\int_Y f(x,y)\,\mathrm{d}\nu(y)\right)\mathrm{d}\mu(x)
=
\int_Y \left(\int_X f(x,y)\,\mathrm{d}\mu(x)\right)\mathrm{d}\nu(y)
=
\int_{X\times Y} f\,\mathrm{d}(\mu\times\nu).
\]
\end{enumerate}
\end{theorem}

\begin{definition}\label{definition:标准磨光核}
设函数 $\rho:\mathbb{R}^n\to\mathbb{R}$ 定义为
\begin{align*}
\rho(x)=
\begin{cases}
c\exp\left(\frac{1}{|x|^2-1}\right), & |x|<1,\\
0, & |x|\geqslant 1,
\end{cases}
\end{align*}
其中$c>0$是使$\int_{\mathbb{R}^n}\rho(x)\,\mathrm{d}x=1$ 的常数. 称 $\rho$ 为 \textbf{标准磨光核}.
\end{definition}

\begin{theorem}[标准磨光核的基本性质]\label{theorem:标准磨光核的基本性质}
对于 $\varepsilon>0$, 定义 $\rho_\varepsilon(x)=\frac{1}{\varepsilon^n}\rho\left(\frac{x}{\varepsilon}\right)$, 其中 $\rho$ 为标准磨光核. 则 $\rho_\varepsilon$ 满足:
\begin{enumerate}[(1)]
\item $\rho_\varepsilon\in C_0^\infty(\mathbb{R}^n)$;
\item $\operatorname{supp}\rho_\varepsilon\subset \overline{B(0,\varepsilon)}$;
\item $\int_{\mathbb{R}^n}\rho_\varepsilon(x)\,\mathrm{d}x=1$;
\item $\rho_\varepsilon\geqslant 0$.
\end{enumerate}
\end{theorem}

\begin{theorem}\label{theorem:实分析--卷积磨光性质}
设 $f\in L_{\mathrm{loc}}^1(\mathbb{R}^n)$, $\varphi\in C_0^\infty(\mathbb{R}^n)$. 则卷积
\begin{align*}
(f*\varphi)(x)\triangleq\int_{\mathbb{R}^n} f(y)\varphi(x-y)\,\mathrm{d}y
\end{align*}
满足 $f*\varphi\in C^\infty(\mathbb{R}^n)$, 且对于任意多重指标 $\alpha\in\mathbb{N}^n$, 有
\begin{align*}
D^\alpha(f*\varphi)=f*(D^\alpha\varphi).
\end{align*}
特别地, 若进一步假设 $f$ 具有紧支集 (即 $f\in L_c^1(\mathbb{R}^n)$), 则 $f*\varphi\in C_0^\infty(\mathbb{R}^n)$.
\end{theorem}
\begin{remark}
卷积满足交换律, 即 $f*\varphi=\varphi*f$. 因此上述光滑性结论对 $\varphi*f$ 同样成立, 且 $D^\alpha(\varphi*f)=(D^\alpha\varphi)*f$ 也成立.
\end{remark}












\end{document}